\section{Evaluasi, Refleksi, dan Integrasi Kompetensi}

Bab akhir ini merupakan titik pengukuran pencapaian kompetensi secara menyeluruh. Sesuai struktur OBE, bab ini berisi kumpulan asesmen (Quiz 1--4, Tugas, UTS, UAS, Project), rubrik penilaian komprehensif, tinjauan pencapaian kompetensi, dan rekomendasi studi lanjutan.

\subsection{Quiz 1 --- Menguji CPMK-1}

Quiz 1 menguji penguasaan \textbf{CPMK-1} (Analisis Kebutuhan Pengguna dan Perancangan Antarmuka), meliputi Sub-CPMK 1.1, 1.2, dan 1.3. Format: \textbf{pilihan ganda}. Materi rujukan: Bab 2--4 (pengantar pemrograman web, kebutuhan fungsional dan non-fungsional, wireframe, mockup, aksesibilitas, prinsip UX).

\subsection{Quiz 2 --- Menguji CPMK-2}

Quiz 2 menguji penguasaan \textbf{CPMK-2} (Implementasi Solusi Web Front-End), meliputi Sub-CPMK 2.1--2.4. Format: \textbf{pilihan ganda}. Materi rujukan: Bab 5--11 (HTML, CSS, form, Flexbox/Grid, JavaScript, DOM, event handling, validasi form, library/framework front-end).

\subsection{Quiz 3 --- Menguji CPMK-3}

Quiz 3 menguji penguasaan \textbf{CPMK-3} (Pengembangan Sistem Back-End dan Integrasi Data), meliputi Sub-CPMK 3.1, 3.2, dan 3.3. Format: \textbf{pilihan ganda}. Materi rujukan: Bab 12--15 (PHP dasar, pemrosesan form, OOP, session/cookie, MySQL, CRUD, integrasi API).

\subsection{Quiz 4 --- Menguji CPMK-4}

Quiz 4 menguji penguasaan \textbf{CPMK-4} (Evaluasi, Deployment, dan Administrasi Web), meliputi Sub-CPMK 4.1, 4.2, dan 4.3. Format: \textbf{pilihan ganda}. Materi rujukan: Bab 15 (pengujian, performa, keamanan XSS/CSRF, deployment, administrasi web).

\subsection{Tugas}

Tugas merupakan asesmen formatif per bab (Tugas 1--4 sesuai RPS), mengacu pada materi yang telah dibahas. Format: \textbf{pilihan ganda}. Tugas mendukung pencapaian Sub-CPMK secara bertahap sesuai bab yang relevan.

\subsection{UTS --- Menguji CPMK-1 dan CPMK-2}

Ujian Tengah Semester (UTS) dilaksanakan pada \textbf{Pertemuan 8}. UTS menguji \textbf{CPMK-1} dan \textbf{CPMK-2} secara komprehensif (analisis kebutuhan, wireframe/mockup, HTML, CSS, JavaScript, library/framework). Format: \textbf{pilihan ganda}. Materi rujukan: Bab 2--11.

\subsection{UAS --- Menguji CPMK-3 dan CPMK-4}

Ujian Akhir Semester (UAS) dilaksanakan pada \textbf{Pertemuan 16}. UAS menguji \textbf{CPMK-3} dan \textbf{CPMK-4} secara komprehensif (PHP, MySQL, API, pengujian, keamanan, deployment). Format: \textbf{pilihan ganda}. Materi rujukan: Bab 12--15.

\subsection{Project --- Menguji CPMK-3}

Project menguji penguasaan \textbf{CPMK-3}. Bentuk: \textbf{tugas membuat aplikasi web sistem informasi akademik (SIA) sederhana}. Aplikasi harus mencakup: back-end PHP, koneksi MySQL, operasi CRUD, session, integrasi API front-end--back-end. Front-end dapat memanfaatkan materi HTML/CSS/JavaScript dari bab sebelumnya. Soal lengkap dan rubrik project tercantum di Lampiran.

\subsection{Rubrik Penilaian Komprehensif Berbasis OBE}

Penilaian mengikuti prinsip OBE: setiap komponen asesmen dipetakan ke CPMK/Sub-CPMK. Quiz, Tugas, UTS, dan UAS (pilihan ganda) dinilai berdasarkan jumlah jawaban benar. Project dinilai berdasarkan rubrik yang terhubung ke Sub-CPMK 3.1, 3.2, 3.3. \textbf{Bobot (sesuai RPS):} Tugas/Praktikum 35\%, Kuis (Quiz 1--4) 10\%, UTS 20\%, Proyek 20\%, UAS 15\%. Rubrik terperinci tercantum di \textbf{Lampiran A} buku ajar.

\subsection{Tinjauan Pencapaian Kompetensi}

Setelah menyelesaikan seluruh bab, mahasiswa diharapkan mampu:

\begin{itemize}
  \item \textbf{CPMK-1}: Mengidentifikasi kebutuhan, membuat wireframe/mockup, menerapkan aksesibilitas dan UX.
  \item \textbf{CPMK-2}: Membangun antarmuka web responsif dan dinamis dengan HTML5, CSS3, JavaScript, dan library/framework.
  \item \textbf{CPMK-3}: Mengembangkan back-end PHP, mengelola MySQL, mengintegrasikan API (diuji oleh Quiz 3, UAS, dan Project).
  \item \textbf{CPMK-4}: Melakukan pengujian, menerapkan keamanan, melakukan deployment dan administrasi web.
\end{itemize}

Tinjau checklist di setiap bab; identifikasi area yang masih perlu dikembangkan; buat rencana belajar lanjutan sebelum UAS dan Project.

\subsection{Rekomendasi Tindak Lanjut Pembelajaran}

Kelulusan mata kuliah Pemrograman Web menandai awal perjalanan --- teknologi web terus berkembang. Rekomendasi studi lanjutan:

\begin{itemize}
  \item \textbf{Framework back-end}: Laravel, Symfony, atau CodeIgniter untuk pengembangan web PHP tingkat lanjut.
  \item \textbf{Framework front-end}: Vue.js, React, atau Angular untuk aplikasi single-page (SPA).
  \item \textbf{DevOps dan CI/CD}: Docker, GitHub Actions, atau Jenkins untuk otomasi deployment.
  \item \textbf{Keamanan web lanjut}: OWASP Top 10 mendalam, penetration testing dasar, keamanan API.
  \item \textbf{Performa dan skala}: caching (Redis, Memcached), load balancing, optimasi database.
  \item \textbf{Referensi standar}: MDN Web Docs, OWASP, web.dev, PHP The Right Way.
\end{itemize}

Teknologi web terus bergerak; seorang developer harus terus belajar dan mengikuti praktik terbaik (best practices).
